% =====================================================================
\section{Literature review and the research gap}
\label{sec:lit}
% =====================================================================

\subsection{Cyclic codes over matrix rings}

The use of a matrix ring as a code alphabet goes back to
Bachoc's construction of modular lattices, and the coding-theoretic study of
cyclic codes over $M_2(\F_2)$ was initiated by Alahmadi, Sboui, Sol\'e and
Yemen \cite{Alahmadi2013CyclicM2F2}, who fixed the convention that cyclic codes
are \emph{right} ideals and used the ring-valued Euclidean dot product. Luo and
Parampalli \cite{LuoParampalli2018M2F2uF2} extended the structure theorem to
$M_2(\F_2+u\F_2)$, i.e. to $t=2$ in our notation, and Dinh, Gaur, Kumar, Singh
and Singh \cite{Dinh2024CyclicRingsOfMatrices} to $M_2(\F_2[u]/\langle u^k\rangle)$
for all $k$ --- the same ring family as the present paper --- obtaining the
structure via an isometry onto $\F_4[u,v]/\langle u^k,v^2\rangle$. Most
recently Joshi, Patel and Prakash \cite{JoshiPatelPrakash2026M4} treated
$M_4(\F_2[u]/\langle u^k\rangle)$ with Euclidean and Hermitian duals and both
the Bachoc and the Gray map into $\F_{16}$. A separate, genuinely different line
uses skew polynomial rings and cyclic algebras over $M_2(\F_2)$
\cite{SiNiu2023SkewM2F2}.

What \emph{none} of these papers contains is a Galois duality, a hull, or a
quantum construction. Their duals are the Euclidean and the Hermitian dual, and
their output is a classical code over a finite field, compared with the best
known classical codes.

\subsection{Duality over finite rings}

The authority for duality over non-commutative finite rings is Wood
\cite{Wood1999Duality}: the extension of the MacWilliams equivalence theorem to
finite rings characterises Frobenius rings, and duality for modules over a
finite Frobenius ring is the correct general framework. We import the statement
that $R=M_m(A_t)$ is Frobenius and prove directly, in
Lemma~\ref{lem:lambda}, that $\Lambda=\lambda\circ\tr$ is a Frobenius functional
on it.

For the weight side, Greferath and Schmidt \cite{GreferathSchmidt1999GrayIsometry}
construct Gray isometries for finite \emph{chain} rings. Their theorem does not
apply to $M_2(A_t)$, which is not a chain ring for $t\ge1$ and $m=2$; we use it
only to delimit what our own map can achieve (Remark~\ref{rem:noisometry}).

\subsection{Galois duality and Galois hulls}

Over a finite field, Fan and Zhang \cite{FanZhang2017GaloisSelfDual} introduced
the $\ell$-Galois inner product $\langle a,b\rangle_\ell=\sum_i a_ib_i^{p^\ell}$
and studied $\ell$-Galois self-dual constacyclic codes; Liu and Pan
\cite{LiuPan2020GaloisHulls} introduced the $\ell$-Galois hull
$\Hull_\ell(C)=C\cap C^{\perp_\ell}$ and gave its dimension formulas. Cao
\cite{Cao2021MDSGaloisHulls} constructed MDS codes with Galois hulls of
arbitrary dimensions and the corresponding entanglement-assisted quantum codes.
All of this takes place over a \emph{commutative} base.

The closest prior work to the present paper is Zhang, Kong and Zheng
\cite{ZhangKongZheng2026QuantumGaloisHulls}, who study constacyclic codes over a
finite non-chain ring $R$ (commutative, with a decomposition
$C=e_1C_1\oplus\cdots\oplus e_tC_t$), determine the dimensions of the
$\ell$-Galois hulls of their Gray images, and construct EAQECCs and QECCs by a
construction $X$ for the Galois inner product, obtaining some QMDS codes and
some QECCs better than the best known ones.

\subsection{The research gap}

The gap is the conjunction of three properties, none of which is present in the
literature:

\begin{enumerate}[label=\textnormal{(\arabic*)},leftmargin=*]
\item the alphabet is \textbf{non-commutative} (a matrix ring over a finite
  chain ring), so that the Galois twist must be an \emph{anti}-automorphism and
  the left/right distinction is substantive;
\item the duality is \textbf{Galois} (a Frobenius parameter $s$ is present,
  interpolating between the Euclidean and Hermitian cases in a way that a
  commutative alphabet cannot reproduce);
\item the outcome is \textbf{quantum}: stabilizer and entanglement-assisted
  codes, obtained from the hull through a Gray map whose induced form is
  identified exactly.
\end{enumerate}

Property (1) alone is the matrix-ring literature; (2) alone is the Galois-hull
literature; (2)+(3) over a commutative non-chain ring is
\cite{ZhangKongZheng2026QuantumGaloisHulls}. The conjunction (1)+(2)+(3) is
what this paper supplies.

Two consequences of the non-commutativity are worth signalling here, because
they have no commutative analogue and they shape the whole paper.

\begin{itemize}
\item The pairing must use $\theta_s=\sigma_s\circ(\cdot)^{\mathsf T}$, not
  $\sigma_s$. Over a commutative ring the transpose is invisible and the
  $\ell$-Galois form is simply Frobenius-in-one-argument; here the transpose is
  what makes $\theta_s$ an anti-automorphism, and the anti-automorphism is what
  makes the dual of a \emph{right} module again a right module
  (Theorem~\ref{thm:duality}). The plain product $\sum_ix_iy_i$, which is the
  Euclidean dot product used in \cite{Alahmadi2013CyclicM2F2}, has a perfect
  dual on left modules but \emph{not} on right modules
  (Lemma~\ref{lem:failure}).
\item The hull need not be principal, so no generator-polynomial formula for
  the hull can exist, and every downstream statement must be proved without a
  principality hypothesis. This is not a technical nuisance: the
  counterexamples are exactly the families that the quantum constructions later
  use (Theorem~\ref{thm:nonprincipal} and Section~\ref{subsec:witnesses}).
\end{itemize}

\subsection{Quantum codes from classical codes}

The stabilizer formalism for $q$-ary quantum codes is due to
\cite{KetkarEtAl2006NonbinaryStabilizer} in the non-binary case. We do not
import the parameter formulas: Section~\ref{sec:quantum} re-derives the
stabilizer lemma and all three routes from the symplectic commutation relation,
so that every hypothesis (in particular the requirement that commutation be
tested on an $\F_p$-basis, Theorem~\ref{thm:fpbasis}) is visible where it is
used.
